\documentclass[12pt]{article}

%----------------- Packages ------------------
\usepackage[utf8]{inputenc}
\usepackage[T1]{fontenc}
\usepackage{amsmath, amssymb}
\usepackage{geometry}
\usepackage{xcolor}
\usepackage{titlesec}
\usepackage{fancyhdr}
\usepackage[breakable]{tcolorbox}
\usepackage{graphicx}
\usepackage{tikz}
\usepackage{float}
\usetikzlibrary{arrows.meta, decorations.markings}

%----------------- Page Setup -----------------
\geometry{a4paper, margin=2.5cm}
\pagestyle{fancy}
\fancyhf{}
\rhead{Final Exam — \textbf{2021}}
\lhead{Analysis V}
\cfoot{\thepage}

%----------------- Title Styling --------------
\titleformat{\section}{\normalfont\Large\bfseries}{Exercise \thesection:}{1em}{}
\titleformat{\subsection}{\normalfont\bfseries}{Answer:}{1em}{}

%----------------- Custom Boxes ----------------
\tcbuselibrary{listingsutf8}
\newtcolorbox{answerbox}{
  colback=gray!10,
  colframe=black,
  fonttitle=\bfseries,
  title=Answer Area,
  breakable,
  before skip=10pt,
  after skip=10pt
}

%----------------- Document Start --------------
\begin{document}

%----------------- Exam Info -------------------
\begin{center}
  \Large\textbf{UNIVERSITY NAME} \\[1em]
  \large\textit{Analysis V — Final Exam} \\[0.5em]
  \large\textit{Year: 2021} \\[2em]
\end{center}

\vspace{0.5cm}

%----------------- EXERCISE 1 ------------------
\section{}
Let \( E = C^1([0,1], \mathbb{R}) \) be the space of continuously differentiable functions and \(\|\cdot\|_\infty\) the uniform convergence norm. Define  

\[
N(f) = |f(0)| + \|f'\|_\infty \quad \text{and} \quad N'(f) = \|f\|_\infty + \|f'\|_\infty.
\]

\begin{enumerate}
    \item Show that \( N \) and \( N' \) are norms on \( E \).
    \item Show that \( N \) and \( N' \) are equivalent.
    \item Are the norms \( N \) and \( N' \) equivalent to the norm \(\|\cdot\|_\infty\)?
\end{enumerate}

\begin{answerbox}
Let $E = C^1([0,1], \mathbb{R})$, $\|\cdot\|_\infty$ the sup norm.

\subsection*{1. Show $N$ and $N'$ are norms}

\textbf{For $N(f) = |f(0)| + \|f'\|_\infty$:}
\begin{itemize}
    \item $N(f) = 0 \Rightarrow |f(0)| = 0$ and $\|f'\|_\infty = 0 \Rightarrow f(0) = 0$ and $f' = 0 \Rightarrow f = 0$.
    \item $N(\lambda f) = |\lambda f(0)| + \|\lambda f'\|_\infty = |\lambda|(|f(0)| + \|f'\|_\infty) = |\lambda| N(f)$.
    \item $N(f+g) = |f(0)+g(0)| + \|f'+g'\|_\infty \leq |f(0)| + |g(0)| + \|f'\|_\infty + \|g'\|_\infty = N(f) + N(g)$.
\end{itemize}
So $N$ is a norm.

\textbf{For $N'(f) = \|f\|_\infty + \|f'\|_\infty$:}
\begin{itemize}
    \item $N'(f) = 0 \Rightarrow \|f\|_\infty = 0$ and $\|f'\|_\infty = 0 \Rightarrow f = 0$.
    \item $N'(\lambda f) = \|\lambda f\|_\infty + \|\lambda f'\|_\infty = |\lambda|(\|f\|_\infty + \|f'\|_\infty) = |\lambda| N'(f)$.
    \item $N'(f+g) \leq \|f+g\|_\infty + \|f'+g'\|_\infty \leq \|f\|_\infty + \|g\|_\infty + \|f'\|_\infty + \|g'\|_\infty = N'(f) + N'(g)$.
\end{itemize}
So $N'$ is a norm.

\subsection*{2. Show $N$ and $N'$ are equivalent}

For $f \in E$, $x \in [0,1]$:
\[
|f(x)| = \left|f(0) + \int_0^x f'(t) dt\right| \leq |f(0)| + \|f'\|_\infty = N(f).
\]
Taking supremum: $\|f\|_\infty \leq N(f)$. Also $|f(0)| \leq \|f\|_\infty$.

Thus:
\[
N(f) = |f(0)| + \|f'\|_\infty \leq \|f\|_\infty + \|f'\|_\infty = N'(f) \leq N(f) + \|f'\|_\infty.
\]
Actually better: $N(f) \leq N'(f)$ and $N'(f) = \|f\|_\infty + \|f'\|_\infty \leq N(f) + N(f) = 2N(f)$.

So $N(f) \leq N'(f) \leq 2N(f)$. Hence $N$ and $N'$ are equivalent.

\subsection*{3. Are $N$, $N'$ equivalent to $\|\cdot\|_\infty$?}

No. Consider $f_n(x) = x^n$ on $[0,1]$.

$\|f_n\|_\infty = 1$ for all $n$, but $N(f_n) = |f_n(0)| + \|f_n'\|_\infty = 0 + n = n \to \infty$.

So $\frac{N(f_n)}{\|f_n\|_\infty} = n \to \infty$, thus $N$ and $\|\cdot\|_\infty$ are not equivalent.

Since $N$ and $N'$ are equivalent, $N'$ is also not equivalent to $\|\cdot\|_\infty$.
\end{answerbox}

%----------------- EXERCISE 2 ------------------
\section{}
Let \((E, \|\cdot\|)\) be a normed vector space.

\begin{enumerate}
    \item Let \((x_n)_{n \in \mathbb{N}}\) be a sequence converging to \( x \) in \( E \) and let \( K \) be a non-empty subset of \( E \) defined by:
    \[
    K = \{x_n \mid n \in \mathbb{N}\} \cup \{x\}.
    \]
    \begin{enumerate}
        \item Show that \( K \) is compact in \( E \) using the definition of compactness.
        \item Show that \( K \) is compact in \( E \) using the Borel property.
    \end{enumerate}
    \item Let \((A_n)_{n \in \mathbb{N}}\) be a decreasing sequence of non-empty compact subsets of \( E \). Show that the intersection \( A = \bigcap_{n \in \mathbb{N}} A_n \) is a non-empty compact subset of \( E \).
\end{enumerate}

\begin{answerbox}
  \begin{enumerate}
    \item $K = \{x_n : n \in \mathbb{N}\} \cup \{x\}$, $x_n \to x$
    \begin{enumerate}
      \item Using definition 
      
      We show every sequence in $K$ has a convergent subsequence in $K$.

      Let $(y_k)$ be a sequence in $K$. If $y_k = x$ infinitely often, we have constant subsequence. Otherwise, $(y_k)$ contains infinitely many $x_n$. Since $x_n \to x$, we can extract subsequence converging to $x \in K$.

      \item Using Borel property 
      
      Let $\{U_i\}_{i \in I}$ be an open cover of $K$. Since $x \in K$, $\exists i_0$ with $x \in U_{i_0}$. As $x_n \to x$, $\exists N$ such that $\forall n \ge N$, $x_n \in U_{i_0}$. The remaining finite set $\{x_1,\dots,x_{N-1}\}$ can be covered by finitely many $U_i$. Thus $K$ is covered by finitely many $U_i$.
      \end{enumerate}
    \item $A_n$ decreasing non-empty compact sets, $A = \bigcap_{n\in\mathbb{N}} A_n$

      Choose $x_n \in A_n$. Since $A_n \subset A_0$ and $A_0$ compact, $(x_n)$ has convergent subsequence $x_{\varphi(n)} \to \ell$.

      For fixed $p$, for $n$ large, $\varphi(n) \ge p$, so $x_{\varphi(n)} \in A_{\varphi(n)} \subset A_p$. Since $A_p$ closed, $\ell \in A_p$. This holds $\forall p$, so $\ell \in \bigcap A_n = A$. Thus $A \neq \emptyset$.

      $A$ is closed (intersection of closed sets) and contained in compact $A_0$, so $A$ is compact.

  \end{enumerate}
\end{answerbox}

%----------------- EXERCISE 3 ------------------
\section{}
Define the function \( f : \mathbb{R}^2 \longrightarrow \mathbb{R} \) by:  

\[
f(x,y) = 
\begin{cases} 
\dfrac{x^2 y^3}{x^2 + y^2} & \text{if } (x,y) \neq (0,0) \\ 
0 & \text{if } (x,y) = (0,0).
\end{cases}
\]

\begin{enumerate}
    \item Compute \(\lim_{(x,y) \to (0,0)} f(x,y)\).
    \item Show that \( f \) is continuous on \( \mathbb{R}^2 \) (justify your answer).
    \item 
    \begin{enumerate}
        \item Determine \( df_{(x,y)} \), the differential of \( f \) for \((x,y) \neq (0,0)\).
        \item Compute \(\frac{\partial f}{\partial x}(0,0)\) and \(\frac{\partial f}{\partial y}(0,0)\).
    \end{enumerate}
    \item Show that \( f \) is of class \( C^1 \) on \( \mathbb{R}^2 \).
    \item Let \( g \) be a function defined on \( \mathbb{R}^2 \) by:
    \[
    g(x,y) = f(x,y) - \frac{1}{2}.
    \]
    \begin{enumerate}
        \item Show that there exists a neighborhood \( U \) of 1 and a \( C^1 \) function \( \mu \) on \( U \) such that:
        \[
        \mu(1) = 1 \quad \text{and} \quad \forall x \in U, \quad g(x,\mu(x)) = 0.
        \]
        \item Compute \( \mu'(1) \), the derivative of \( \mu \) at the point 1.
    \end{enumerate}
\end{enumerate}

\begin{answerbox}
\begin{enumerate}
  \item Limit at $(0,0)$
For $(x,y) \neq (0,0)$:
\[
|f(x,y)| = \frac{x^2 |y|^3}{x^2 + y^2} = x^2 |y| \cdot \frac{y^2}{x^2 + y^2} \le x^2 |y| \cdot 1 \le (x^2 + y^2)^{3/2}.
\]
As $(x,y) \to (0,0)$, $|f(x,y)| \to 0$. So $\lim_{(x,y)\to(0,0)} f(x,y) = 0 = f(0,0)$.

  \item Continuity on $\mathbb{R}^2$

$f$ is rational on $\mathbb{R}^2\setminus\{(0,0)\}$, so continuous there. From part 1, continuous at $(0,0)$. Thus continuous everywhere.

  \item Differential
  \begin{enumerate}
    \item For $(x,y) \neq (0,0)$:
\[
\frac{\partial f}{\partial x} = \frac{2x y^3(x^2+y^2) - x^2 y^3 \cdot 2x}{(x^2+y^2)^2} = \frac{2x y^5}{(x^2+y^2)^2},
\]
\[
\frac{\partial f}{\partial y} = \frac{3x^2 y^2(x^2+y^2) - x^2 y^3 \cdot 2y}{(x^2+y^2)^2} = \frac{3x^4 y^2 + x^2 y^4}{(x^2+y^2)^2}.
\]
Thus $df_{(x,y)}(h,k) = \frac{\partial f}{\partial x}h + \frac{\partial f}{\partial y}k$.

    \item At $(0,0)$:
\[
\frac{\partial f}{\partial x}(0,0) = \lim_{t\to0} \frac{f(t,0)-f(0,0)}{t} = \lim_{t\to0} \frac{0}{t} = 0.
\]
\[
\frac{\partial f}{\partial y}(0,0) = \lim_{t\to0} \frac{f(0,t)-f(0,0)}{t} = \lim_{t\to0} \frac{0}{t} = 0.
\]
  \end{enumerate}
  \item $C^1$ on $\mathbb{R}^2$

Need to show partial derivatives are continuous at $(0,0)$.

For $\frac{\partial f}{\partial x}$:
\[
\left|\frac{2x y^5}{(x^2+y^2)^2}\right| \le 2|x||y| \cdot \frac{y^4}{(x^2+y^2)^2} \le 2|x||y| \to 0.
\]
So $\lim_{(x,y)\to(0,0)} \frac{\partial f}{\partial x}(x,y) = 0 = \frac{\partial f}{\partial x}(0,0)$.

For $\frac{\partial f}{\partial y}$:
\[
\left|\frac{3x^4 y^2 + x^2 y^4}{(x^2+y^2)^2}\right| \le \frac{3x^4 y^2}{(x^2+y^2)^2} + \frac{x^2 y^4}{(x^2+y^2)^2} \le 3y^2 + x^2 \to 0.
\]
So $\lim_{(x,y)\to(0,0)} \frac{\partial f}{\partial y}(x,y) = 0 = \frac{\partial f}{\partial y}(0,0)$.

Thus $f \in C^1(\mathbb{R}^2)$.

  \item Implicit function
  \begin{enumerate}
    \item $g(x,y) = f(x,y) - \frac{1}{2}$

At $(1,1)$: $f(1,1) = \frac{1^2 \cdot 1^3}{1^2+1^2} = \frac{1}{2}$, so $g(1,1) = 0.
$$\frac{\partial g}{\partial y}(1,1) = \frac{\partial f}{\partial y}(1,1) = \frac{3\cdot1^4\cdot1^2 + 1^2\cdot1^4}{(1^2+1^2)^2} = \frac{3+1}{4} = 1 \neq 0$$
$
By implicit function theorem, $\exists U \ni 1$, $C^1$ function $\mu: U \to \mathbb{R}$ with $\mu(1)=1$ and $g(x,\mu(x))=0$ for $x\in U$.

    \item Compute $\mu'(1)$

From $g(x,\mu(x)) = 0$, differentiate:
\[
\frac{\partial f}{\partial x}(1,1) + \frac{\partial f}{\partial y}(1,1) \cdot \mu'(1) = 0.
\]
Now $\frac{\partial f}{\partial x}(1,1) = \frac{2\cdot1\cdot1^5}{(1^2+1^2)^2} = \frac{2}{4} = \frac{1}{2}$.

So $\frac{1}{2} + 1 \cdot \mu'(1) = 0 \Rightarrow \mu'(1) = -\frac{1}{2}$.
    \end{enumerate}
\end{enumerate}
\end{answerbox}

%----------------- END ------------------
\end{document}